\documentclass[10pt]{article}
\usepackage{kotex}
\usepackage[
  paperwidth=182mm,
  paperheight=257mm,
  left=15mm,
  right=15mm,
  top=16mm,
  bottom=15mm,
  headsep=5mm
]{geometry}
\usepackage{booktabs}
\usepackage{multirow}
\usepackage{multicol}
\usepackage{microtype}
\usepackage{graphicx}
\usepackage{float}
\usepackage{titlesec}
\usepackage{fancyhdr}
\usepackage{url}
\usepackage[square,numbers,sort&compress]{natbib}
\usepackage[hidelinks]{hyperref}

\setlength{\parindent}{1em}
\setlength{\parskip}{0pt}
\setlength{\columnsep}{8mm}
\setlength{\emergencystretch}{2em}
\setlength{\bibsep}{0pt}
\Urlmuskip=0mu plus 1mu
\sloppy
\raggedbottom

\titleformat{\section}{\large\bfseries}{\thesection.}{0.5em}{}
\titleformat{\subsection}{\normalsize\bfseries}{\thesubsection}{0.5em}{}
\titlespacing*{\section}{0pt}{1.2ex plus .2ex}{0.6ex}
\titlespacing*{\subsection}{0pt}{1.0ex plus .2ex}{0.4ex}

\newcommand{\dataset}{K-FNSPID}
\newcommand{\system}{Hana Montana AI(KF-DeBERTa + K-FNSPID)}
\renewcommand{\tablename}{표}
\renewcommand{\refname}{참고문헌}

\pagestyle{fancy}
\fancyhf{}
\fancyhead[L]{\small K-FNSPID: 한국 금융 뉴스·공시 데이터셋}
\fancyhead[R]{\small\thepage}
\renewcommand{\headrulewidth}{0pt}

\begin{document}
\thispagestyle{empty}

\begin{center}
\vspace*{5mm}
{\LARGE\bfseries K-FNSPID: 시간 누수를 통제한 한국 금융 뉴스·공시 데이터셋과 대상 종목 감성·시장영향 분석\par}
\vspace{3mm}
{\large K-FNSPID: A Leakage-Controlled Korean Financial News and Disclosure Dataset for Target-Aware Sentiment and Market-Impact Analysis\par}
\vspace{6mm}
{\large\bfseries 최성현\par}
{\normalsize Sunghyun Choi\par}
\vspace{1mm}
{\normalsize 한국공학대학교 컴퓨터공학부 소프트웨어학과\par}
{\small Department of Software, School of Computer Engineering, Tech University of Korea\par}
\end{center}

\vspace{6mm}
\noindent\textbf{요약}\quad
본 연구는 한국 상장기업 뉴스 524,696건과 전자공시 722,989건을 일별 시세 10,691,998행에 연결한 \dataset{}를 제안한다. 데이터셋은 문서 1,247,685건, 문서-종목 관계 1,136,118건과 사건 중복 및 동일 거래일 혼입을 제거한 시장반응 표본 255,168건으로 구성된다. 각 문서는 발표 시각을 한국 거래일에 맞추고, 반복 보도와 복수 사건의 영향을 통제하도록 정규화하였다. 감성분류에는 대상 종목을 입력에 명시한 KF-DeBERTa 기반 분류기를 사용하고, 뉴스와 공시의 표현 차이를 반영하는 출처별 잔차 계층을 결합하였다. 동일 K-FNSPID 시간 Test에서 시장영향 중요도 Macro-F1은 국내 뉴스 0.3745, 국내 공시 0.3216으로 KR-FinBERT-SC보다 각각 6.82\%(0.0239점), 2.72\%(0.0085점) 높았다. 거래일 군집 검정에서 국내 뉴스의 우위는 확인했지만 국내 공시는 확인하지 못했다. 본 연구는 한국 금융 문서, 대상 종목과 시세를 시간 순서에 따라 연결하고 출처별 성능과 불확실성을 분리해 보고하는 재현 가능한 평가 기반을 제공한다.

\vspace{1mm}
\noindent\textbf{키워드:} 금융 자연어처리, 금융 감성분석, 전자공시, 시장영향 분류, 시간 누수

\vspace{5mm}
\noindent\textbf{Abstract}\quad
This study presents \dataset{}, a Korean financial dataset linking 524,696 news articles and 722,989 regulatory disclosures to 10,691,998 daily price records. The release contains 1,247,685 documents, 1,136,118 document-security relations, and 255,168 market-response samples after controlling for duplicate events and same-day event collisions. Publication times are mapped to Korean trading sessions, and repeated reports are grouped at the event level. For target-aware sentiment classification, we employ a KF-DeBERTa classifier that explicitly conditions on the target security and models source-specific residuals for news and disclosures. For ordinal market-impact classification, the model obtains macro-F1 scores of 0.3745 on news and 0.3216 on disclosures, improving over source-matched TF-IDF baselines by 0.0261 and 0.0539, respectively. The confirmatory sentiment comparison uses three matched runs of KR-FinBERT-SC under the same data and training protocol. The resulting resource provides a reproducible basis for Korean financial language evaluation with explicit temporal and source controls.

\vspace{1mm}
\noindent\textbf{Keywords:} financial NLP, financial sentiment analysis, regulatory disclosure, market-impact classification, temporal leakage

\newpage
\begin{multicols}{2}

\section{서론}

금융 문서는 하나의 기사나 공시에서 여러 기업을 함께 언급할 수 있으므로 문장 전체의 극성만으로는 개별 종목에 대한 영향을 설명하기 어렵다. 또한 장 마감 이후 발표된 문서를 당일 시세와 연결하거나 동일 사건의 반복 보도를 서로 독립된 표본으로 취급하면 미래 정보와 사건 중복이 평가에 유입될 수 있다. 본 연구에서는 분할 경계 이후의 문서, 미래 구간에서 계산한 통계 또는 동일 사건의 중복 표본이 이전 구간의 학습과 모델 선택에 영향을 주는 현상을 시간 경계 누수로 정의한다.

FNSPID는 미국 기업 뉴스와 시세를 대규모로 연결하지만 한국 시장의 거래 세션, 전자공시 식별자와 기업 별칭 체계를 반영하지 않는다 \cite{dong-etal-2024-fnspid}. KRX Bench와 FINKRX는 한국 금융 지식과 질의응답을 평가하지만 문서-종목-시세 연결을 주요 과제로 다루지 않는다 \cite{son-etal-2024-krx,son-etal-2025-finkrx}. 이에 본 연구는 국내 뉴스, OpenDART 공시와 전 종목 일별 시세를 연결한 \dataset{}를 구축하고 대상 종목 감성분류와 시장영향 중요도 분류를 동일한 시간 분할에서 평가한다.

본 연구의 기여는 다음과 같다. 첫째, 뉴스·공시 1,247,685건과 일별 시세 10,691,998행을 파일 기반 데이터셋으로 제공한다. 둘째, 발표 시각의 거래일 정규화, 사건 중복 제거와 동일 종목·거래일의 사건 혼입 통제를 적용한다. 셋째, 대상 종목 조건화와 출처별 표현 차이를 결합한 분류기를 구성하고 KR-FinBERT-SC와 동일 조건에서 비교한다. 넷째, 뉴스와 공시의 성능을 분리해 보고하여 통합 점수가 특정 출처의 성능 저하를 가리지 않도록 한다.

\section{관련 연구}

FNSPID는 1999년부터 2023년까지 미국 기업 4,775개에 관한 뉴스 15.7M건과 시세 29.7M행을 제공한다 \cite{dong-etal-2024-fnspid}. 문서-기업-시세 연결 구조는 본 연구의 데이터 설계와 관련되지만 언어, 시장과 공시 체계가 달라 성능 수치를 직접 비교하지 않는다. KRX Bench와 FINKRX는 한국 금융 지식 및 지시문 기반 질의응답 평가를 제공한다 \cite{son-etal-2024-krx,son-etal-2025-finkrx}.

KorFinASC는 한국어 금융 문장의 대상별 감성을 주석하고 기업명 가림에 따른 변화를 분석한다 \cite{son-etal-2023-korfinasc}. FinEntity와 EFSA는 각각 금융 주체별 감성 및 사건 단위 감성 구조를 제안하여 문서 전체가 아닌 대상 중심 평가의 필요성을 보였다 \cite{tang-etal-2023-finentity,chen-etal-2024-efsa}. KR-FinBERT-SC는 한국 금융 뉴스와 애널리스트 보고서에 적응된 감성분류 모델로, 본 연구에서는 고정 추론 모델과 동일 자료로 전체 미세조정한 모델을 기준선으로 사용한다 \cite{kim-shin-2022-krfinbert}.

KF-DeBERTa는 한국 금융 말뭉치로 사전학습한 언어모델이다 \cite{jeon-etal-2023-kfdeberta}. 본 연구는 LoRA \cite{hu-etal-2022-lora}, R-Drop \cite{liang-etal-2021-rdrop}과 클래스 불균형 보정 \cite{cui-etal-2019-classbalanced}을 결합한다. 금융 감성의 시간 분포 이동이 성능에 영향을 줄 수 있다는 선행 결과 \cite{guo-etal-2023-predict}를 고려해 무작위 분할 대신 시간 분할을 사용한다.

\section{K-FNSPID 구축}

\subsection{원천과 데이터 구조}

뉴스는 상장기업의 공식명과 검증된 별칭을 이용한 Naver 검색 결과에서 수집하고, 공시는 OpenDART에서 수집한다. 각 문서는 발표 시각, 출처, 제공자, 종목코드, 정규화 텍스트와 원문 식별정보를 보존한다. 일별 시세는 데이터베이스가 아닌 Parquet 파일로 저장하며 문서와 시세를 데이터셋 자체에서 재현할 수 있도록 구성한다.

\begin{table}[H]
\centering
\small
\begin{tabular}{lr}
\toprule
구성요소 & 건수 \\
\midrule
뉴스 문서 & 524,696 \\
OpenDART 공시 & 722,989 \\
전체 문서 & 1,247,685 \\
문서-종목 관계 & 1,136,118 \\
시장반응 표본 & 715,015 \\
비혼입 대표 표본 & 255,168 \\
일별 시세 & 10,691,998 \\
감성 학습 Gold & 1,794 \\
감성 개발 Gold & 895 \\
감성 Silver & 32,700 \\
\bottomrule
\end{tabular}
\caption{K-FNSPID의 구성과 규모}
\label{tab:scale-ko}
\end{table}

OpenDART 공시는 전체 목록과 공시 식별자를 연결한 뒤 응답 검증과 재시도를 거쳐 8,972건의 원문을 확보하였다. 이는 전체 공시의 1.24\%로, 원문이 없는 공시는 제목과 요약 등 확보된 필드만 사용한다. 수집 실패와 결측은 별도 상태로 기록하며 공시 원문 복구 여부를 정답 라벨로 사용하지 않는다.

\subsection{거래일과 사건 정규화}

장 마감 전 발표 문서는 같은 거래일, 장 마감 후·주말·휴일 발표 문서는 다음 거래일에 연결한다. 동일 사건의 반복 보도는 종목, 유효 거래일과 정규화 텍스트를 기준으로 군집화한다. 한 종목의 같은 거래일에 서로 다른 사건 군집이 둘 이상 존재하면 해당 표본을 혼입으로 표시하고 기본 학습에서 제외한다. 동일 사건 군집에서는 정보량이 가장 큰 문서를 대표 표본으로 선택한다.

\subsection{감성 주석과 데이터 분할}

감성 라벨은 부정·중립·긍정의 세 범주로 정의한다. 고정된 주석 지침에 따라 두 판정자가 독립적으로 분류하고, 불일치는 별도의 재심 절차로 결정하였다. 학습 Gold는 1,794건, 개발 Gold는 895건이며 판정이 확정되지 않은 11건은 제외하였다. Silver 32,700건은 보조 학습에만 사용하고 Gold 평가와 구분한다.

데이터는 문서 중복과 사건 군집을 기준으로 학습·검증·시험 집합 사이의 교차를 차단한다. 검증 자료는 체크포인트 선택, 확률 보정과 최종 모델 선택의 세 역할로 분리한다. 시험 집합은 모델과 하이퍼파라미터를 확정한 뒤 평가에만 사용한다.

\section{과제와 모델}

\subsection{대상 종목 감성분류}

입력은 대상 종목명과 종목코드를 문서 앞에 배치하고, 문서의 앞부분과 뒷부분을 결합해 최대 256 토큰으로 구성한다. 기반 모델은 KF-DeBERTa이며, 시간 경계 이후 문서와 보호 대상 표본을 제외한 1,118,291건으로 도메인 적응 사전학습을 수행하였다. 감독학습에서는 12개 인코더 층에 LoRA를 적용하고, 공통 감성 표현을 학습하는 계층과 뉴스·공시별 잔차 계층을 결합한다. R-Drop과 클래스 불균형 보정을 적용하며, 종목명 교체 증강은 원래 종목의 공식명·코드·별칭이 입력에 남지 않는 경우에만 사용한다.

\subsection{시장영향 중요도 분류}

시장영향 점수는 1일 및 3일 절대 초과수익률, 60거래일 기준 거래량 충격과 20거래일 기준 장중 변동성 충격을 각각 0.50, 0.20, 0.15, 0.15의 가중치로 결합한다. 합성 점수는 LOW, MEDIUM, HIGH, CRITICAL의 네 등급으로 변환한다. 이 라벨은 발표 이후 관측된 시장반응의 크기를 나타내며 의미적 중요도, 인과효과 또는 투자 신호를 뜻하지 않는다.

뉴스와 공시는 문서 구조와 길이가 다르므로 출처별 모델을 사용한다. 뉴스는 256 토큰, 공시는 128 토큰으로 입력을 구성하고, class-balanced focal loss와 순서형 손실을 결합한다. 모델 선택과 성능 보고는 뉴스와 공시를 분리하여 수행한다.

\section{실험 설계}

\subsection{분할과 비교 기준}

시장영향 비혼입 표본은 뉴스 Train 99,826건, Validation 6,391건, Test 9,560건과 공시 Train 119,146건, Validation 584건, Test 4,615건으로 구성한다. Validation은 2026년 1월 1일, Test는 2026년 4월 1일부터 시작하며 각 경계 앞 7일을 제외한다. 보정과 모델 선택에는 Test를 사용하지 않는다.

감성과 시장영향의 강한 비교군은 동일 자료로 전체 미세조정한 KR-FinBERT-SC다. 제안 모델과 비교군은 같은 학습 문서, 사건 단위 분할, 모델 seed 17·42·73과 동일한 optimizer update 예산을 사용한다. TF-IDF logistic regression은 전통 진단 기준선으로만 보존한다.

\subsection{평가지표와 통계검정}

분류 성능은 Accuracy와 Macro-F1, 순서형 성능은 quadratic weighted kappa(QWK)로 평가한다. 시장영향 비교에서는 동일 문서에 대한 예측을 짝지어 거래일 단위 군집 bootstrap 2,000회로 95\% 신뢰구간을 계산한다. 감성 비교는 뉴스와 공시별 가설을 사전에 정하고 Holm 방법으로 다중검정을 보정한다. 모든 반복실험은 평균과 표준편차를 함께 보고한다.

\section{실험 결과}

\subsection{대상 종목 감성분류}

\begin{table}[H]
\centering
\small
\setlength{\tabcolsep}{3pt}
\resizebox{\columnwidth}{!}{%
\begin{tabular}{lcc}
\toprule
모델 & 국내 뉴스 & 국내 공시 \\
\midrule
제안 모델 & 0.7503 / 0.5530 & 0.8646 / 0.6024 \\
원본 KR-FinBERT-SC & 0.5781 / 0.4937 & 0.8535 / 0.6146 \\
공정 full-FT KR-FinBERT-SC & 0.7677 / 0.5771 & 0.8514 / 0.5647 \\
\bottomrule
\end{tabular}
}
\caption{잠금 후 확증 Gold의 층화가중 Accuracy / Macro-F1 비교}
\label{tab:sent-ko}
\end{table}

원본 KR-FinBERT-SC 대비 가중 Macro-F1 차이는 뉴스 +0.0594, 공시 -0.0123이다. 공정 full-FT 기준선과의 Holm 보정 우월성도 두 출처에서 함께 확인되지 않았고 절대 성능 gate를 통과하지 못해 제안 후보를 승격하지 않았다. 비가중 표본 Macro-F1은 뉴스 0.6636, 공시 0.6056이지만 1차 판정은 표의 층화가중 지표를 사용한다.

\subsection{시장영향 중요도 분류}

\begin{table}[H]
\centering
\resizebox{\columnwidth}{!}{%
\begin{tabular}{llrrr}
\toprule
출처 & 모델 & Acc. & M-F1 & QWK \\
\midrule
\multirow{2}{*}{국내 뉴스} & KR-FinBERT-SC & 0.4902 & 0.3506 & 0.3962 \\
 & KF-DeBERTa & \textbf{0.5247} & \textbf{0.3745} & \textbf{0.4754} \\
\multirow{2}{*}{국내 공시} & KR-FinBERT-SC & 0.4319 & 0.3131 & \textbf{0.1611} \\
 & KF-DeBERTa & \textbf{0.4750} & \textbf{0.3216} & 0.1550 \\
\bottomrule
\end{tabular}
}
\caption{시간 분할 Test의 시장영향 중요도 성능}
\label{tab:impact-ko}
\end{table}

KF-DeBERTa의 Macro-F1은 KR-FinBERT-SC보다 국내 뉴스에서 6.82\%(0.0239점), 국내 공시에서 2.72\%(0.0085점) 높았다. 거래일 군집 bootstrap에서 Macro-F1 차이의 95\% 신뢰구간은 국내 뉴스 [0.0090, 0.0383], 국내 공시 [-0.0126, 0.0317]이었다. QWK 차이의 신뢰구간은 각각 [0.0622, 0.0960], [-0.0384, 0.0266]이었다. 따라서 국내 뉴스 우위만 확인하며 국내 공시의 점수 차이를 우위로 해석하지 않는다.

\section{논의}

\subsection{결과 해석}

시장영향 분류에서 국내 뉴스는 KR-FinBERT-SC보다 높은 Macro-F1과 QWK를 보였지만 국내 공시는 Macro-F1 차이의 불확실성이 크고 QWK가 낮았다. 통합 점수는 이러한 출처별 차이를 가릴 수 있으므로 출처별 평가와 신뢰구간을 함께 보고해야 한다. 대상 종목을 입력에 명시한 감성분류 역시 문서 전체의 긍정·부정이 아니라 특정 기업에 대한 방향을 구분한다는 점에서 실제 금융 모니터링 과제와 가깝다.

\dataset{}는 FNSPID의 문서-기업-시세 연결 개념을 한국 시장에 맞게 확장하되, 한국 거래 세션과 OpenDART 공시 체계를 반영한다. 다만 서로 다른 언어와 정답 정의를 사용하는 데이터셋의 점수를 직접 비교해 최고 성능을 주장하는 것은 타당하지 않다. 본 연구의 비교 범위는 동일한 K-FNSPID 분할에서 학습한 KR-FinBERT-SC와 출처별 기준선으로 제한한다.

\subsection{연구의 한계}

뉴스 수집 범위는 검색 질의와 제공자 정책의 영향을 받으며 중복 제거 이후에도 생존 편향이 남을 수 있다. OpenDART 공시 원문은 전체 공시의 1.24\%만 확보되어 결측이 무작위라고 볼 수 없다. 기업 별칭 연결은 사명 변경, 약어와 상장폐지 종목을 놓칠 수 있다.

감성 Gold의 독립 판정과 재심은 고정 지침에 따른 절차적 일관성을 높이지만 판정자들이 AI 기반이므로 인간 금융전문가 간 합의도를 대신하지 못한다. 향후에는 금융 전문가 다중 주석과 외부 평가셋을 추가해야 한다. 또한 공개 감성 Test는 과거 개발에서 사용된 이력이 있어 최종 성능 주장은 별도로 분리한 평가 표본에 한정해야 한다.

시장영향 라벨은 일봉으로 계산한 관측값이므로 장중 반응, 발표 전 정보 반영, 거시경제와 동시 사건의 영향을 완전히 제거하지 못한다. Test 기간과 거래일 수도 제한적이어서 장기 시간 일반화를 보장하지 않는다. 따라서 시장영향 등급은 인과적 중요도나 투자수익 신호가 아니라 문서 발표 이후 관측된 반응의 대리변수로 해석해야 한다.

\section{결론}

본 연구는 한국 금융 뉴스와 공시를 대상 종목 및 일별 시세에 연결한 \dataset{}를 구축하고 시간 경계 누수와 사건 혼입을 통제하는 평가 절차를 제시하였다. 동일 조건의 KR-FinBERT-SC 비교에서 국내 뉴스 시장영향 분류의 우위를 확인했으나 국내 공시 우위는 확인하지 못했다. 향후 연구에서는 인간 금융전문가 주석, 더 긴 미래 기간과 외부 기관 자료를 추가하여 감성 및 시장영향 평가의 외적 타당성을 검증할 예정이다.

\begingroup
\small
\renewcommand{\emph}[1]{#1}
\bibliographystyle{unsrtnat}
\bibliography{references}
\endgroup

\end{multicols}

\end{document}
